\section{Introduction}\label{sec:introduction}

This section establishes the research problem and derives the research
questions that guide this study. We begin by situating Generative AI (GenAI)
in professional practice, then examine its implications for Enterprise
Architecture (EA) and Reference Architecture (RA) development. Two gaps
emerge: a practice gap, where no empirical research documents how EA
practitioners integrate GenAI into RA development, and a guidance vacuum,
where practitioners experiment in isolation without shared frameworks.

The \textit{Theoretical Background} (Section~\ref{sec:theoretical-background}) provides the theoretical foundation
on EA, RAs, GenAI capabilities, and distributed cognition.
The \textit{Conceptual Model} (Section~\ref{sec:conceptual-model}) presents the conceptual model that frames
the investigation. The \textit{Methodology} (Section~\ref{sec:methodology}) explains the Design-Oriented
Delphi methodology used to answer the research questions. The following
subsections detail this investigation:

\localtableofcontents\clearpage

\subsection{The emergence of Generative AI in Professional Practice}

Since the public release of ChatGPT in November 2022, Generative AI (GenAI) has
become a routine part of professional work
\parencite{brynjolfsson_generative_2025}. ChatGPT reached 100 million users
within two months \parencite{hu_chatgpt_2023}. By 2025, 84\% of developers use
AI coding tools, up from 76\% one year earlier
\parencite{stack_overflow_2024, stack_overflow_2025}, and 75\% of knowledge
workers across 31 countries report using AI regularly
\parencite{microsoft_wti_2025}. Yet adoption has outpaced preparation: 83\% of
workers are self-taught in AI, and only 52\% of senior leaders have deployed
formal AI training \parencite{ey_agentic_2025}.

The infrastructure behind this adoption is large and growing. In 2025, the major
US hyperscalers (Amazon, Microsoft, Google, and Meta) are collectively investing
approximately \$320 billion in AI infrastructure, more than 1\% of US GDP
\parencite{understanding_ai_2025}. Microsoft alone has committed \$80 billion to
data center construction in fiscal year 2025 \parencite{cnbc_tech_2025}. The
\textcite{iea_energy_2025} projects that global electricity consumption from
data centers will more than double by 2030, reaching 945 TWh (equivalent to
Japan's entire current consumption). In the United States, data centers are
forecast to account for almost half of electricity demand growth between 2024
and 2030 \parencite{iea_energy_2025}. \textcite{mckinsey_cost_2025} estimates
the total cost to scale AI data center infrastructure at \$7 trillion. These
investments have multi-decade lifespans. The practical question is not whether
GenAI will persist, but how practitioners should adapt their work to it.

\subsection{Reference Architecture Development}

Enterprise Architecture (EA) is a discipline that aligns organizational
strategy, processes, and IT infrastructure
\parencite{lankhorst_enterprise_2017}. EA practitioners analyze, design, and
govern how organizations structure their technology ecosystems
\parencite{the_open_group_togaf_2022}. Among their outputs are Reference
Architectures (RAs): documented patterns and templates for system
implementations \parencite{cloutier_concept_2010}. RAs enable reuse across
projects, reduce design risk, promote interoperability, and serve as knowledge
repositories for stakeholder communication
\parencite{cloutier_concept_2010, angelov_framework_2012}.

In practice, RA development involves more than formal EA titles.
\textcite{kotusev_enterprise_2018} showed that EA artifacts are built
collaboratively by architects, business executives, analysts, and developers.
Solution architects, platform engineers, and senior developers routinely shape
these artifacts through daily work. Developing RAs requires analyzing existing
systems, identifying patterns, engaging stakeholders, and synthesizing technical
and business considerations.

GenAI is already entering this work. Developer tools powered by GenAI are
shaping architectural patterns from the bottom up, through code generation,
design suggestions, and automated documentation, often before formal EA
governance becomes aware
\parencite{peng_impact_2023, tabarsi_llms_reshaping_se_2025}. This makes
understanding how practitioners use GenAI in RA development urgent.

\subsection{The Practice Gap}

No empirical research documents how EA practitioners integrate GenAI into RA
development\footnote{ EBSCO host advanced search conducted on June 8, 2025,
    resulted in 0 publications using the query: ((enterprise architecture OR ea
    practice) AND (gen ai OR generative artificial intelligence OR llm OR large
    language models)) AND SO (MIS Quarterly Executive OR Information Systems
    Research OR Journal of the Association for Information Systems OR European
    Journal of Information Systems) Filters applied: peer-reviewed articles,
    publication date from January 1, 2020, to June 8, 2025
}. Research in adjacent fields (software engineering, management consulting,
general knowledge work) shows that GenAI changes professional practice
\parencite{peng_impact_2023, dellacqua_navigating_2023, noy_experimental_2023},
but EA remains unstudied
\parencite{banh_affordance_genai_ea_2025, ettinger_ea_dynamic_capability_2025}.

This gap matters for three reasons. First, EA work involves both technical and
social dimensions. Practitioners navigate organizational politics, synthesize
diverse perspectives, and create artifacts that must resonate with multiple
stakeholders. How GenAI tools, primarily trained on technical content, interact
with this sociotechnical reality is unknown. Second, RA development is
knowledge-intensive. It requires capturing, synthesizing, and codifying
architectural knowledge. GenAI's capabilities for pattern recognition and
content generation could change how this knowledge work is done, but no evidence
documents whether or how. Third, research in other domains shows that
professionals restructure their work practices around AI rather than simply
using it as a faster tool \parencite{noy_experimental_2023,
tankelevitch_metacognitive_2024, dellacqua_navigating_2023}. Whether EA
practitioners are experiencing similar changes is unexamined.

\subsection{The Guidance Vacuum}

Practitioners who experiment with GenAI in EA work do so without shared
vocabulary, evaluation criteria, or validated guidance. Specifically, they lack:

\begin{enumerate}
    \item Evidence of which practices work for GenAI use in EA
          \parencite{banh_affordance_genai_ea_2025}.
    \item Frameworks for judging when GenAI helps versus hinders architectural
          activities \parencite{esposito_genai_architecture_2026}.
    \item Guidance on managing the risks of human-AI collaboration in EA,
          particularly over-reliance \parencite{dellacqua_navigating_2023}.
\end{enumerate}

Without shared guidance, each practitioner reinvents the wheel
\parencite{feldman_theorizing_2011}. Emerging practices cannot consolidate into
reliable routines, and organizations risk both missed opportunities and quality
problems \parencite{dellacqua_navigating_2023}.

\subsection{Research Opportunity}

In our view, these gaps call for an investigation into how EA practitioners
develop new practices as they integrate GenAI into RA development. This is not
a tool-adoption study; Section~\ref{subsec:framework-dynamics} develops this argument further. Research in adjacent domains shows that professionals
reshape their work around AI capabilities rather than bolting AI onto existing
workflows \parencite{noy_experimental_2023, dellacqua_navigating_2023}.

We focus on RA development because it is a central EA activity
\parencite{cloutier_concept_2010, the_open_group_togaf_2022} that spans the
full range of architectural work: analysis, synthesis, and documentation. If
GenAI changes how RAs are developed, the effects will extend to governance,
implementation, and knowledge management. Beyond documenting current practices,
the aim is to develop practical guidance that supports EA practitioners in
navigating this change.

In sum, GenAI is a structural shift in professional knowledge work. RA
development is knowledge-intensive sociotechnical work for which no empirical
guidance exists. Developer-level GenAI use is already shaping architectural
artifacts before formal EA governance becomes aware. The following subsections
translate this problem into research questions.

\subsection{Deriving the Research Questions}\label{sec:research-question}

The introduction identified two gaps. A \emph{practice gap}: no empirical
research documents how EA practitioners integrate GenAI into RA development.
And a \emph{guidance vacuum}: practitioners experiment in isolation, without
shared frameworks or vocabulary. These gaps call for research that both
documents emerging practices and produces actionable guidance.

\subsection{Main Research Question}

\begin{quote}
\textit{What practices do EA practitioners adopt when using GenAI for RA development,
and how can these be systematized into guidance?}
\end{quote}

The first part of this question targets practice discovery. The second targets
guidance development. This reflects a practice-based perspective: technology
integration is not a one-time event but an ongoing process shaped by how people
actually use tools in context \parencite{feldman_theorizing_2011,
orlikowski_using_2000}. GenAI and EA practice shape one another. Practitioners
adapt their workflows to GenAI, while GenAI outputs depend on practitioner
expertise and organizational context \parencite{orlikowski_sociomaterial_2007}.
Answering this question requires a methodology that captures both current
practices and the ambition to produce practical guidance.

\subsection{Sub-Questions}\label{subsec:sub-questions}

Four sub-questions decompose the main RQ into a progression from description
through analysis and interpretation to design:

\vspace{0.5em}
\noindent\textbf{sRQ1:} What GenAI practices are EA practitioners currently using when developing RAs?

\vspace{0.5em}
\noindent\textbf{sRQ2:} What changes to EA practice are EA practitioners experiencing when RAs are developed using GenAI?

\vspace{0.5em}
\noindent\textbf{sRQ3:} How do EA practitioners make sense of impacts and trade-offs when adopting GenAI for RA development?

\vspace{0.5em}
\noindent\textbf{sRQ4:} How should effective GenAI practices for RA development be systematized into guidance for EA practitioners?

\vspace{0.5em}
\noindent Throughout these questions, ``EA practitioners'' refers broadly to all
professionals contributing to RA development, as formally defined in
Section~\ref{sec:theoretical-background}.

\subsection{Sub-Question Progression}

The four sub-questions follow a deliberate sequence. sRQ1 is
\emph{descriptive}: it maps which GenAI practices EA practitioners currently
use. sRQ2 is \emph{analytical}: it examines how these practices change existing
EA work. sRQ3 is \emph{interpretive}: it surfaces how practitioners weigh the
impacts and trade-offs they experience.

sRQ4 differs from the first three. Where sRQ1 through sRQ3 investigate what
\emph{is}, sRQ4 asks what \emph{should be}. It synthesizes empirical findings
into a guidance artifact that practitioners can use
\parencite{gregor_positioning_2013}. This shift from empirical inquiry to
artifact construction is characteristic of the Design-Oriented Delphi
methodology adopted in this study (Section~\ref{sec:methodology}).
